\section{Criterios de efectividad para un dato de descenso}
\label{VIII.7}

Consideremos, como de costumbre, un morfismo de preesquemas
$$
g\colon S'\to S
$$
y un $S'$-preesquema~$X'$. De acuerdo con los hechos
\ifthenelse{\boolean{orig}}
{generales~VII,~9,}
{generales~VII,~9\footnote{Este número no existe; véase la exposición 212
del Seminario Bourbaki (1962)},}
dar un dato de descenso sobre~$X'$ con respecto a~$g$ equivale a dar un par
de equivalencia \cite{VIII.3}:
$$
\xymatrix@C=.5cm{ q_{1},q_{2}\colon X''\ar@/^2pt/[r]\ar@/_2pt/[r] & X'}
$$
tal que el morfismo estructural~$X'\to S'$ sea compatible con este par y con
el par de equivalencia
$$
\xymatrix@C=.5cm{p_{1},p_{2}\colon
S''=S'\times_{S}S'\ar@/^2pt/[r]\ar@/_2pt/[r] & S'}
$$
definido por~$g$, y tal que los dos cuadrados (o uno de ellos, lo que viene
a ser lo mismo por simetría) extraídos del diagrama correspondiente
$$
\xymatrix{X'\ar[d] & X''\ar@/^2pt/[l]\ar@/_2pt/[l]\ar[d]\\ S' &
S''\ar@/^2pt/[l]\ar@/_2pt/[l]}
$$
(utilizando bien~$p_{1},q_{1}$, bien~$p_{2},q_{2}$)
\ifthenelse{\boolean{orig}}
{sean \emph{cartesianos}}
{sean \emph{cartesianos}}.
Una solución del problema de descenso planteado por este
dato de descenso, \ie un objeto~$X$ sobre~$S$ provisto de un
isomorfismo~$X\times_{S}S'\isomfrom X'$ compatible con los datos de descenso,
equivale a dar un cuadrado \emph{cartesiano}
$$
\xymatrix{X\ar[d] & X'\ar[l]_{h}\ar[d]\\ S & S'\ar[l]_{g}}
$$
que satisface~$hq_{1}=hq_{2}$.

Como
\marginpar{220}
la clase de los morfismos fielmente planos y cuasicompactos es estable por
cambio de base, y un morfismo fielmente plano y cuasicompacto es un
epimorfismo efectivo en virtud de~\Ref{VIII.5.3}, se sigue de la teoría
general \cite{VIII.D}:

\begin{proposition}
\label{VIII.7.1}
Supongamos que~$g\colon S'\to S$ es fielmente plano y cuasicompacto. Para que
un dato de descenso sobre~$X'$ con respecto a~$g$ sea efectivo, es necesario
y suficiente que la relación de equivalencia~$R=(q_{1},q_{2})$ que este define
sea efectiva (\ie que exista el cociente~$X'/R$ y que~$X''$ se convierta en
el cuadrado fibrado de~$X'$ sobre~$X'/R$), y que el morfismo
canónico~$X'\to X'/R$ sea fielmente plano y cuasicompacto.
\end{proposition}

Así, la cuestión de la efectividad de un dato de descenso es un caso
particular de la cuestión de la efectividad de un grafo de equivalencia, y
por esta vía pueden obtenerse diversos criterios de efectividad dados en
este número. Sin embargo, en el contexto del descenso se dispone del
teorema~\Ref{VIII.2.1}, que implica que \emph{si $X'$ es afín sobre~$S'$,
todo dato de descenso sobre~$X'$ con respecto a~$g$ es efectivo}; este
enunciado no tiene análogo para el paso al cociente por un grafo de
equivalencia plano general. Todos los criterios de efectividad que damos
aquí pueden considerarse asimismo deducidos del anterior.

Sea~$U'$ un subpreesquema de~$X'$ (o, más generalmente, un subobjeto de~$X'$
en la categoría~$\Sch$); se dice que~$U'$ es \emph{estable por el dato de
descenso} sobre~$X'$ si puede dotarse a~$U'$ de un dato de descenso con
respecto a~$g$ tal que la inmersión~$U'\to X'$ sea compatible con los datos
de descenso. Esto significa también que las imágenes inversas de~$U'$ en
~$X''$ por~$q_{1}$ y~$q_{2}$ son iguales (o también, como se dice, que~$U'$
es \emph{estable por la relación de equivalencia}~$R$); desde luego, el dato
de descenso en cuestión sobre~$U'$ es entonces único y se llama \emph{dato
de descenso inducido}
\index{dato de descenso inducido|hyperpage}%
por el de~$X'$. Sentado esto:

\begin{proposition}
\label{VIII.7.2}
Sea~$(X_{i}')$ un recubrimiento de~$X'$ por abiertos~$X_{i}'$ estables
\marginpar{221}
por el dato de descenso. Para que el dato de descenso sobre~$X'$ sea
efectivo, es necesario y suficiente que lo mismo suceda con los datos de
descenso inducidos en los~$X_{i}'$.
\end{proposition}

Esta es, por ejemplo, una consecuencia fácil de~\Ref{VIII.7.1}; se deja al
lector el detalle de la demostración.

\begin{corollaire}
\label{VIII.7.3}
Sea~$(S_{i})$ un recubrimiento abierto de~$S$, y para todo~$i$ sean~$S_{i}'$
y~$X_{i}'$ deducidos de~$S'$
\ifthenelse{\boolean{orig}}
{de}
{y}
$X'$ mediante el cambio de base~$S_{i}\to S$. Para que el dato de descenso
sobre~$X'$ sea efectivo, es necesario y suficiente que, para todo~$i$, el
dato de descenso sobre~$X_{i}'$ con respecto a
~$g_{i}\colon S_{i}'\to S_{i}$ sea efectivo.
\end{corollaire}

Este criterio nos reduce siempre, en la práctica, al caso en que~$S$ es
afín. En el caso en que~$S'$ también es afín, que es el más frecuente en las
aplicaciones, se tiene:

\begin{corollaire}
\label{VIII.7.4}
Supongamos que~$S$ y~$S'$ son afines. Para que el dato de descenso sobre~$X'$
sea efectivo, es necesario y suficiente que~$X'$ sea reunión de
abiertos~$X_{i}'$ afines y estables por el dato de descenso.
\end{corollaire}

La suficiencia proviene de~\Ref{VIII.7.2} y del hecho de que, si~$X_{i}'$ es
afín, es afín sobre~$S'$ y puede aplicarse~\Ref{VIII.2.1}. Para la necesidad,
obsérvese que, si~$X'$ proviene de~$X$ y~$X$ está recubierto por abiertos
afines~$X_{i}$, entonces los~$X_{i}'=X_{i}\times_{S}S'$ son abiertos afines
estables por los datos de descenso y recubren~$X'$.

\begin{corollaire}
\label{VIII.7.5}
Sea~$g\colon S'\to S$ un morfismo fielmente plano, cuasicompacto y
\emph{radicial}. Entonces~$g$ es un morfismo de descenso \emph{efectivo},
\ie para todo~$X'$ sobre~$S'$, todo dato de descenso sobre~$X'$ con respecto
a~$g\colon S'\to S$ es efectivo.
\end{corollaire}

En efecto, en virtud de~\Ref{VIII.7.3} puede suponerse que~$S$ es afín; como
~$S'$ es radicial sobre~$S$ y, por tanto, separado, $S'$ es separado. Por
otra parte, para todo~$x'\in X'$, la fibra
~$R(x')=q_{2}(q_{1}^{-1}(x'))$ de la relación de equivalencia conjuntista
definida
\marginpar{222}
por la relación de equivalencia~$R$ se reduce a un punto. En efecto, como~$g$
es radicial, lo mismo sucede con~$p_{1},p_{2}$, que se deducen de aquel por
el cambio de base~$S'\to S$, y, por tanto, también con~$q_{1},q_{2}$, que se
deducen de los anteriores por el cambio de base~$X'\to S''$. Así,
\emph{todo abierto de~$X'$ es estable} por el dato de descenso. Recubramos
entonces~$X'$ por abiertos afines~$X_{i}'$; son afines
\ifthenelse{\boolean{orig}}
{sobre~$S$}
{sobre~$S'$}
puesto que~$S'$ es separado, por lo que el dato de descenso inducido es
efectivo según~\Ref{VIII.2.1}. Se concluye entonces mediante~\Ref{VIII.7.2}.

Obsérvese que~\Ref{VIII.7.5} da el único caso conocido de un morfismo de
descenso efectivo en la categoría de los preesquemas, y probablemente sea
en efecto el único caso, incluso limitándose a los esquemas noetherianos o a
los esquemas de tipo finito sobre un cuerpo.

Cuando se supone~$S$ localmente noetheriano y~$S'$ de tipo finito sobre~$S$,
el enunciado~\Ref{VIII.7.5} es también un caso particular del siguiente (que
generaliza el descenso galoisiano de Weil y el descenso inseparable de
Cartier):

\begin{corollaire}
\label{VIII.7.6}
Sea~$g\colon S'\to S$ un morfismo finito localmente libre (\ie definido por
un álgebra sobre~$S$ que es un módulo localmente libre de tipo finito) y
sobreyectivo (por tanto, $g$ es fielmente plano y cuasicompacto, luego es un
morfismo de descenso). Sea~$X'$ un~$S'$-preesquema provisto de un dato de
descenso. Para que este dato sea efectivo, es necesario y suficiente que,
para todo~$x'\in X'$, la fibra~$R(x')=q_{2}(q_{1}^{-1}(x'))$ esté contenida
en un abierto afín (condición que se satisface automáticamente si~$X'$ es
cuasiproyectivo sobre~$S'$).
\end{corollaire}

La observación entre paréntesis proviene del hecho de que, si~$s$ es el
punto de~$S$ que es imagen de~$x'$, entonces~$R(x')$ es finito y está
contenido en la
\ifthenelse{\boolean{orig}}
{fibra~$X_{s}$,}
{fibra~$X'_{s}$,}
por otra parte, como~$X'$ es cuasiproyectivo sobre~$S'$ y~$S'$ es finito
sobre~$S$, $X'$ es cuasiproyectivo sobre~$S$, lo cual implica que una fibra
\ifthenelse{\boolean{orig}}
{de~$X/S$}
{de~$X'$ sobre~$S$}
está contenida en un abierto afín.

Como todo subconjunto finito de un esquema afín admite un sistema
fundamental de entornos afines, se ve que no se pierde la hipótesis al
restringirse sobre un abierto afín de~$S$; esto, en virtud
de~\Ref{VIII.7.3}, nos reduce al caso en que~$S$ es afín. En virtud
de~\Ref{VIII.7.4}, queda por demostrar
\ifthenelse{\boolean{orig}}
{que~$X'$}
{que~$x'$}
está contenido
\marginpar{223}
en un abierto afín \emph{estable} por el dato de descenso. Sea, en efecto,
~$U$ un abierto afín que contiene~$R(x')$; entonces el saturado
$$
\ifthenelse{\boolean{orig}}
{R(X-U)=q_{2}(q_{1}^{-1}(X-U))}
{R(X'-U)=q_{2}(q_{1}^{-1}(X'-U))}
$$
no corta a~$R(x')$. Por otra parte, como~$q_{2}$ es finito (pues~$g$ y, por
tanto, $p_{2}$ lo son) y, por consiguiente, cerrado, el segundo miembro es
un subconjunto cerrado de~$X'$. Sea~$U'$ su complementario en~$X'$; es,
pues, un abierto \emph{saturado}, y se tiene
$$
R(x')\subset U'\subset U
$$
con~$U$ afín, aunque~$U'$ no es afín \emph{a priori}. Como el subconjunto
finito~$R(x')$ del esquema afín~$U$ posee un sistema fundamental de entornos
afines de la forma~$U_{f}$, se ve, reemplazando~$f$ por su restricción a
~$U'$, que existe una sección~$f$ de~$\mathcal{O}_{U}$ tal que:
$$
R(x')\subset U_{f}',\quad U_{f}'\textrm{ es afín}.
$$
Sea entonces~$U''=q_{1}^{-1}(U')=q_{2}^{-1}(U')$; denotemos todavía
por~$q_{1},q_{2}$ los morfismos inducidos~$U''\to U'$ y consideremos
$$
f'=\Norm_{q_{2}}(q_{1}^*(f))\quoi,
$$
donde~$\Norm_{q_{2}}$ designa la \emph{norma} con respecto al morfismo
finito localmente libre~$q_{2}\colon U''\to U'$. La compatibilidad de la
formación de la norma con el cambio de base implica fácilmente que~$f'$ es
una sección \emph{invariante}:
$$
q_{1}^*(f')=q_{2}^*(f')
$$
lo cual implica que~$U_{f'}'$ es un abierto saturado de~$U'$. Más
precisamente, si~$Z(f')$ designa el conjunto de ceros de una sección~$f'$,
se obtiene, en virtud de las propiedades de las normas:
$$
Z(f')=q_{2}(Z(q_{1}^*(f)))=q_{2}(q_{1}^*(Z(f)))=
R(U'-U_{f}')\quoi,
$$
lo cual
\marginpar{224}
implica que~$U_{f'}'=U'-Z(f')$ es saturado, contiene~$R(x')$ y está
contenido en~$U_{f}'$. Como este último es afín, se sigue que~$U_{f'}'$ lo
es también (pues es igual a~$(U_{f}')_{f''}$, con
~$f''=f'|U_{f'}'$). Es, por tanto, un abierto afín saturado que contiene
~$R(x')$ y, por consiguiente, a~$x'$, lo que concluye la demostración.
\par
